\section*{Chapitre \ref{ch:qm:estimation} --- Estimation}

\begin{solution}{exo:qm:estimation:1}
$\hat\lambda = 100/50 = 2$ par seconde~; $\mathrm{se} = \hat\lambda/\sqrt{100} = 0.2$.
\end{solution}

\begin{solution}{exo:qm:estimation:2}
Par jour $\mathrm{SR}_1 = 1/\sqrt{252} = 0.063$ et $\mathrm{se} = \sqrt{(1 + 0.002)/1\,008} = 0.0315$~; annualisé, $0.0315\sqrt{252} = 0.50$.
\end{solution}

\begin{solution}{exo:qm:estimation:3}
Biaisé~: sa moyenne vaut $(n - 1)\sigma^2/n$. Convergent~: le biais et la variance tendent vers zéro lorsque $n \to \infty$.
\end{solution}

\begin{solution}{exo:qm:estimation:4}
La variance de long terme vaut $\sigma^2(1 + \rho)/(1 - \rho) = 1.5\sigma^2$, si bien que l'\omterm{def:qm:estimation:estimator}{erreur type} naïve est $\sqrt{1.5} = 1.22$ fois trop petite.
\end{solution}

\begin{solution}{exo:qm:estimation:5}
$\mathcal I(p) = 1/(p(1 - p))$, donc $\mathrm{se} = \sqrt{0.55 \times 0.45/200} = 0.035$ et $t = 0.05/0.035 = 1.42$~: non significatif.
\end{solution}

\begin{solution}{exo:qm:estimation:6}
$\Var(\hat\sigma^2) \approx 2\sigma^4/n$ et $\ln\hat\sigma = \tfrac12\ln\hat\sigma^2$, donc $\Var(\ln\hat\sigma) \approx \tfrac14 \cdot 2\sigma^4/(n\sigma^4) = 1/(2n)$~: $\mathrm{se} =
1/\sqrt{4\,000} = 0.0158$, la valeur hessienne du chapitre, correcte pour des données gaussiennes uniquement.
\end{solution}

\begin{solution}{exo:qm:estimation:7}
62.5~\% avec l'\omterm{def:qm:estimation:estimator}{erreur type} iid, 92.2~\% avec Newey--West à sept retards et 93.5~\% à vingt.
\end{solution}

\begin{solution}{exo:qm:estimation:8}
Deux rendements glissants sur 20 jours consécutifs partagent 19 jours~: la série est fortement autocorrélée, si bien que son écart type sous-estime
l'incertitude de sa moyenne, et il n'y a qu'environ $3 \times 252/20 = 38$ observations indépendantes. Annualiser un rendement sur 20 jours
par $\sqrt{252}$ gonfle aussi le \omterm{def:qm:estimation:sharpe}{ratio de Sharpe} d'un facteur $\sqrt{20}$~: utiliser des rendements sur 20 jours sans chevauchement annualisés par $\sqrt{252/20}$, ou le P\&L quotidien,
avec une \omterm{def:qm:estimation:estimator}{erreur type} HAC.
\end{solution}

\begin{solution}{pb:qm:estimation:1}
\textbf{1.} Moyenne 4.03 pb, écart type 42.3 pb.
\textbf{2.} $42.3/\sqrt{1\,260} = 1.19$ pb~; $t = 3.38$.
\textbf{3.} $\widehat{\mathrm{SR}} = 1.51$, \omterm{def:qm:estimation:estimator}{erreur type} iid 0.45.
\textbf{4.} 0.80, 0.59, 0.37, 0.16 contre 0.8, 0.6, 0.4, 0.2.
\textbf{5.} Des jours consécutifs conservent quatre cinquièmes des mêmes positions.
\textbf{6.} $1 + 2(0.8 + 0.6 + 0.4 + 0.2) = 5$~; \omterm{def:qm:estimation:estimator}{erreur type} vraie $42\sqrt{5/1\,260} = 2.65$ pb.
\textbf{7.} $\lfloor4 \times 12.6^{2/9}\rfloor = 7$ retards~: 2.32 pb et $t = 1.73$.
\textbf{8.} 2.33 pb~: les pondérations de Bartlett au-delà du quatrième retard n'ajoutent presque rien.
\textbf{9.} 0.87, environ le double de la valeur iid.
\textbf{10.} 62.5~\% du temps.
\textbf{11.} $\hat\sigma\sqrt{5/n} = 42.3\sqrt{5/1\,260} = 2.66$ pb.
\textbf{12.} $4/2.65 = 1.51$.
\textbf{13.} $n = (3 \times 42\sqrt5/4)^2 = 4\,960$ jours, soit environ 20 ans.
\textbf{14.} Les autocovariances empiriques sont biaisées vers zéro dans les échantillons finis, et les pondérations de Bartlett réduisent encore celles qui
comptent le plus (retards 1 à 4).
\textbf{15.} Les rendements sans chevauchement sont indépendants, si bien que la formule iid est correcte pour eux, au prix d'une faible perte d'efficacité~;
l'estimation de la moyenne reste presque inchangée.
\textbf{16.} Regroupement de volatilité (hétéroscédasticité), changements de régime et tendances lentes dans l'edge du signal.
\textbf{17.} Le seuil de significativité~: vingt variantes testées rendent un $t$ de 1.7, voire de 3.4, bien moins surprenant (chapitre~12).
\textbf{18.} La moyenne, son \omterm{def:qm:estimation:estimator}{erreur type} HAC et sa largeur de bande, la statistique $t$, le \omterm{def:qm:estimation:sharpe}{ratio de Sharpe} avec son \omterm{def:qm:estimation:estimator}{erreur type} robuste,
et le nombre de variantes testées.
\textbf{19.} Résultat nommé~: \emph{la statistique $t$ qui a été divisée par deux}~: la statistique $t$ iid de 3.38 devient 1.73 avec les \omterm{def:qm:estimation:estimator}{erreurs type}
de Newey--West (et 1.51 avec la vraie variance de long terme), car des positions qui se chevauchent rendent la variance de long terme cinq
fois la variance.
\textbf{20.} Lorsque les observations sont dépendantes ou non identiquement distribuées, comme avec des positions qui se chevauchent.
\end{solution}

\begin{solution}{iq:qm:estimation:1}
Choisir le paramètre qui rend les données observées les plus probables. Il est convergent, asymptotiquement gaussien et efficace (il
atteint la borne de Cramér--Rao), invariant par reparamétrisation, et fournit les \omterm{def:qm:estimation:estimator}{erreurs type} à partir de la courbure de la
log-vraisemblance.
\iqlookfor{la définition et la propriété d'efficacité.}
\end{solution}

\begin{solution}{iq:qm:estimation:2}
Avec des rendements quotidiens indépendants, son \omterm{def:qm:estimation:estimator}{erreur type} vaut environ $1/\sqrt{\text{années}} = 1$~: l'estimation se situe à une \omterm{def:qm:estimation:estimator}{erreur type} de
zéro, soit une probabilité unilatérale d'environ 84~\% que la vraie valeur soit positive sous un a priori plat, loin d'une preuve.
\iqlookfor{$\mathrm{se} \approx 1/\sqrt T$ pour les \omterm{def:qm:estimation:sharpe}{ratios de Sharpe} annuels.}
\end{solution}

\begin{solution}{iq:qm:estimation:3}
$A^{-1}BA^{-1}$, avec $A$ la hessienne espérée de l'objectif et $B$ la variance de son gradient. Lorsque la vraisemblance est
mal spécifiée, ou avec les moindres carrés en présence d'hétéroscédasticité, $A \ne B$ et l'expression usuelle $A^{-1}$ est fausse~; la \omterm{def:qm:estimation:sandwich}{variance sandwich}, elle, ne l'est pas.
\iqlookfor{pain et garniture, et l'égalité de l'information.}
\end{solution}

\begin{solution}{iq:qm:estimation:4}
Des fenêtres qui se chevauchent rendent la variable dépendante autocorrélée jusqu'à la longueur de la fenêtre, si bien que les résidus le sont aussi~: les
\omterm{def:qm:estimation:estimator}{erreurs type} usuelles sont trop petites, d'un facteur voisin de la racine carrée du recouvrement. Utiliser Newey--West avec au moins autant de retards que le
recouvrement (ou Hansen--Hodrick), ou des observations sans chevauchement.
\iqlookfor{autocorrélation induite par le recouvrement et HAC.}
\end{solution}

\begin{solution}{iq:qm:estimation:5}
La variance du score, égale à l'opposé de la hessienne espérée de la log-vraisemblance~: la courbure, c'est-à-dire l'information
qu'une observation apporte sur $\theta$. Cramér--Rao~: la variance d'un \omterm{def:qm:estimation:estimator}{estimateur sans biais} est au moins $\mathcal I^{-1}/n$.
\iqlookfor{la courbure égale l'information.}
\end{solution}

\begin{solution}{iq:qm:estimation:6}
Le paramètre $\theta^*$ qui minimise la \omterm{def:qm:estimation:kl}{divergence de Kullback--Leibler} de la distribution vraie vers le modèle~: la meilleure
approximation au sein du modèle. Ses \omterm{def:qm:estimation:estimator}{erreurs type} sont celles de la \omterm{def:qm:estimation:sandwich}{variance sandwich}, et non l'inverse de la hessienne.
\iqlookfor{pseudo-valeur vraie et variance robuste.}
\end{solution}
